
\documentclass[twocolumn]{article}
\usepackage{amsmath, amssymb, amsthm, graphicx}
\usepackage{geometry}
\geometry{margin=0.75in}

\title{\textbf{Quantifying Critical Transition Regions via Curvature-Dependent Width Scaling in Complexity Functions}}
\author{}
\date{}

\begin{document}
\maketitle

\begin{abstract}
We introduce a formal framework to quantify critical transition regions in systems described by a scalar complexity function $C(x)$. We derive a curvature-dependent scaling law linking the width of the transition region to local geometric properties.
\end{abstract}

\section{Introduction}
Understanding transitions in complex systems is central across physics, information theory, and nonlinear dynamics.

\section{Theory}
Let $C: I \subset \mathbb{R} \to \mathbb{R}$ be twice differentiable.

\section{Main Result}
\begin{equation}
W_\epsilon \approx 2\sqrt{\frac{2\epsilon C_{\max}}{|C''(x^*)|}}
\end{equation}

\section{Multidimensional Extension}
\begin{equation}
(x-x^*)^T Q (x-x^*) \le 2\epsilon C_{\max}
\end{equation}

\section{Conclusion}
We generalize criticality from points to regions and derive a universal scaling law.

\section{References}
Shannon (1948), Landau \& Lifshitz (1980), Strogatz (2018)

\end{document}
